%% THERMODYNAMIQUE %%
%
\section[Temp\'erature de l'air \`a partir de 
la temp\'erature �sonique� et du rapport de m�lange]
{Temp\'erature de l'air \`a partir de 
la temp\'erature sonique et du rapport de m�lange en vapeur d'eau}

\itemstyle{ Op\'erateur~: } {\sf Tair\_soni}
\index{Tair\_soni}

%%%%%OBJET%%%%%
\itemstyle{ Objet~: }
Calcul de la temp\'erature de l'air, \`a partir de la
temp\'erature sonique, et du rapport de m�lange en vapeur d'eau.

%%%%%ARGUMENTS%%%%%
\bigskip
%%%%%ENTREES%%%%
\begin{tabarg}
\multicolumn{3}{l}{\titretab{Entr\'ees}}\\
\hline
$T_m$ & Vecteur, {\tt float} & Temp\'erature sonique mesur\'ee ($K$ ou $^\circ C$) \\
$r$ & Vecteur, {\tt float} & Rapport de m\'elange en vapeur d'eau ($g/kg$)\\
\end{tabarg}

\medskip
%%%%%SORTIES%%%%
\begin{tabarg}
\multicolumn{3}{l}{\titretab{Sorties}}\\
\hline
$T$ & Vecteur, {\tt float} & Temp\'erature de l'air (m\^eme unit\'e que $T_m$)\\
\end{tabarg}


%%%%%FORMULE%%%%
\itemstyle{ Formule }

%%% Temperature de l'air \`a partir de la tempe sonique

On calcule d'abord une humidit� sp�cifique en 
\textit{kg}/\textit{kg} :
$ q= \frac{10^{-3}r}{1+10^{-3}r}$

Puis : \[ T = (1 - 0.518 . q).T_m \] 
(Avec, si n�cessaire, conversion 
des degr�s Celsius vers Kelvin et retour).


%%%%%COMMENTAIRE%%%%
\itemstyle{Commentaire~:}
 On appelle temp�rature "sonique" une temp�rature calcul�e
� partir de la vitesse de propagation du son dans l'air
en assimilant l'air � un gaz parfait sec : $c = \sqrt{ \gamma_a R_a T}$
donne $T = \frac{c^2}{\gamma_a R_a}$.

%%%%%AUTEUR%%%%%
\itemstyle{ Auteur~:}
CNRM/GMEI/TRAMM

%%%%%REFERENCES%%%%%
\itemstyle{ R\'ef\'erences~:}
Triplet-Roche \cite{Triplet-Roche}, Kaimal 1990 \cite{Kaimal-1990}, 
Pr\'esentation \`a l'atelier intrumentation de 1996 \cite{AIE96-TAA}.
